%%=============================================================================
%%    Praktische Uitwerking
%%=============================================================================

%    Steunend op de inleiding en de conceptuele oplossing bespreek je nu stap voor stap de uitwerking van je
%    project, waarbij je in detail zal gaan om bepaalde zaken toe te lichten. Een grondige uitleg waarom je
%    bepaalde stappen op die manier aanpakt is essentieel. Minstens even belangrijk is een kritische interpretatie
%    van de resultaten die elke stap oplevert.

%    De praktische uitwerking kan omvatten:
%    ● Uitleg van de werking van een programma of website
%    ● Ontwerp van een elektronische schakeling in detail (berekeningen, printontwerp, test, …)
%    ● De essentiële stukken van de ontwikkeling het programma (leg dit vooral uit a.d.h.v. flowcharts, 
%       beperk de bespreking van zuivere programmacode tot een minimum of tot de relevante of vernieuwende delen)
%    ● Bespreking van meet- of testopstellingen, met enkele essentiële uitgevoerde metingen
%    ● Bespreking van een zelf opgesteld en uitgevoerd stappenplan voor onderhoud of installatie.
%    Merk op:
%    ● Het is belangrijk de opgetreden problemen en daarbij horende oplossingen te vermelden.
%    ● Verwerk waar nodig het cijfermateriaal statistisch en/of breng het passend in beeld.
%    ● De opsplitsing naar voorstudie en praktische uitwerking kun je eventueel per deelopdracht
%       behandelen.

%    Sluit steeds af met een overzicht van wat je gedaan hebt in dit hoofdstuk en wat het resultaat was.
%    Meer informatie is te vinden in het hoofdstuk “Specifieke hoofdstukken” van “Rapporteren voor technici” pg.
%    90 tot 95

\chapter{\IfLanguageName{dutch}{Praktische Uitwerking}{}}%
\label{ch:praktische-uitwerking}

Steunend op de inleiding en de conceptuele oplossing bespreek je nu stap voor stap de uitwerking van je
project, waarbij je in detail zal gaan om bepaalde zaken toe te lichten. Een grondige uitleg waarom je
bepaalde stappen op die manier aanpakt is essentieel. Minstens even belangrijk is een kritische interpretatie
van de resultaten die elke stap oplevert.

De praktische uitwerking kan omvatten:
\begin{itemize}[]
    \item Uitleg van de werking van een programma of website
    \item Ontwerp van een elektronische schakeling in detail (berekeningen, printontwerp, test, …)
    \item De essentiële stukken van de ontwikkeling het programma (leg dit vooral uit a.d.h.v. flowcharts, beperk de bespreking van zuivere programmacode tot een minimum of tot de relevante of vernieuwende delen)
    \item Bespreking van meet- of testopstellingen, met enkele essentiële uitgevoerde metingen
    \item Bespreking van een zelf opgesteld en uitgevoerd stappenplan voor onderhoud of installatie.
\end{itemize}

Merk op:
\begin{itemize}[]
    \item Het is belangrijk de opgetreden problemen en daarbij horende oplossingen te vermelden.
    \item Verwerk waar nodig het cijfermateriaal statistisch en/of breng het passend in beeld.
    \item De opsplitsing naar voorstudie en praktische uitwerking kun je eventueel per deelopdracht behandelen.
\end{itemize}

Sluit steeds af met een overzicht van wat je gedaan hebt in dit hoofdstuk en wat het resultaat was.
Meer informatie is te vinden in het hoofdstuk “Specifieke hoofdstukken” van “Rapporteren voor technici” pg.
90 tot 95